% ===================================================================
%  TABLICA Nr 16 — Wielki magazyn. — Bazar. — Zabawki.
%  Meme regime que les precedents. Tableau a UNE SEULE COLONNE, avec
%  une note en pied placee entre le 8 et le 9.
%
%  LE BLOC t16-15-1 EST L'UN DES TROIS ECARTS DECLARES, et pour une
%  raison de COMPOSITION : le fac-simile ido y ecrit « d) » sans sa
%  parenthese ouvrante, si bien que le releve des renvois ne le voit
%  pas. Les traductions ecrivent le renvoi entier — le neerlandais et
%  le tcheque le font — et cette colonne aussi.
%
%  CE MEME BLOC EST LE SEUL DU LIVRET A NUMEROTER PAR LETTRES : de
%  (a) a (f), pour la table d'astronomie. Les lettres se cliquent sur
%  la page de lecture comme les numeros.
%
%  ET LE 5 EST UN COURS DE STRATEGIE EN TROIS ALINEAS ou presque
%  chaque mot est en gras sans porter de renvoi : siege, bombardement,
%  assaut, sortie, retraite, armistice. Ce ne sont pas des termes de
%  planche, ce sont les mots que le grand-pere apprend a l'enfant. On
%  garde le gras : il est dans le fac-simile.
% ===================================================================

%%K t16-apar-1 apar
\VUsaut{4.0mm}
\VUtitre{15.0pt}{60}{TABLICA Nr 16}
\VUsaut{2.0mm}
\VUpk{13.2pt}{40}{Wielki magazyn. --- Bazar.}
\VUpk{13.2pt}{40}{Zabawki.}
\VUsaut{3.4mm}

%%K t16-01-1 p
1. --- Nowy rok się zbliża. Wielkie magazyny przygotowały swoje wystawy \VUgras{zabawek}
\textsuperscript{(1)} i podarków noworocznych.

%%K t16-01-2 p
Poprosiłem mojego dziadka, żeby zaprowadził mnie do bazaru zobaczyć tę piękną wystawę, i
zgodził się na to.

%%K t16-02-1 p
2. --- Najpierw zatrzymaliśmy się przed pięknymi tarczami z bronią, które mieszczą
\VUgras{karabin}, \VUgras{kepi}, \VUgras{szablę}, \VUgras{pas do szpady} i parę
\VUgras{epoletów}, albo \VUgras{hełm}, \VUgras{pancerz} \textsuperscript{(3)} i
\VUgras{pistolet} \textsuperscript{(4)}. To wszystko zajęło mnie o wiele bardziej niż
\VUgras{latarnie czarnoksięskie} \textsuperscript{(2)}.

%%K t16-tit-1 sub
\VUsaut{3.0mm}
\VUpk{10.2pt}{40}{Piękne widoki.}
\VUsaut{2.6mm}

%%K t16-03-1 p
3. --- W tym samym dziale był \VUgras{wartownik} \textsuperscript{(5)} w swojej
\VUgras{budce}; zdawał się pełnić wartę przed całą menażerią z tektury. Ileż zwierząt
tam było: \VUgras{Papuga} \textsuperscript{(6)} na swojej grzędzie, \VUgras{nosorożec}
\textsuperscript{(7)} z rogiem na nosie, \VUgras{renifer} \textsuperscript{(8)},
\VUgras{struś} \textsuperscript{(9)}, \VUgras{małpa} \textsuperscript{(10)} zgryzająca
orzech, \VUgras{lis} \textsuperscript{(11)} niosący kurę, \VUgras{krokodyl}
\textsuperscript{(12)} rozwierający ogromne szczęki, \VUgras{wielbłąd}
\textsuperscript{(13)}, elegancki \VUgras{chart} \textsuperscript{(14)},
\VUgras{pantera} \textsuperscript{(15)}, potem dwa zwierzęta, których nie znałem, a
które są, jak mówią, \VUgras{łasicą} \textsuperscript{(16)} i \VUgras{hieną}
\textsuperscript{(17)}. Także tam były \VUgras{lampart} \textsuperscript{(18)} i
\VUgras{wilk} \textsuperscript{(21)}; zupełnie blisko były milutkie \VUgras{szczurki}
\textsuperscript{(19)} i maleńkie \VUgras{myszki} \textsuperscript{(20)} kauczukowe. Na
koniec \VUgras{hipopotam} \textsuperscript{(22)} i garbaty \VUgras{dromader}
\textsuperscript{(23)} zamykały zbiór. Zostałbym tam jeszcze wiele godzin, gdyby obok
nie było wspaniałego obozu pełnego \VUgras{ołowianych żołnierzyków}
\textsuperscript{(29)}, który przyciągnął moją uwagę.

%%K t16-tit-2 sub
\VUsaut{4.0mm}
\VUpk{10.2pt}{40}{Żołnierze i wojna.}
\VUsaut{2.8mm}

%%K t16-04-1 p
4. --- Mój dziadek, który jest \VUgras{inwalidą} \textsuperscript{(24)} i który musiał
opuścić armię z powodu \VUgras{rany,} za którą otrzymał \VUgras{medal wojskowy,} bardzo
lubi opowiadać mi o \VUgras{kampaniach,} w których brał udział. Był szczęśliwy,
tłumacząc rozmaite ruchy małej armii w miniaturze. Grupa prawdziwych żołnierzy
zwiedzających bazar: \VUgras{saper} \textsuperscript{(25)}, \VUgras{strzelec alpejski}
\textsuperscript{(26)} z głową nakrytą \VUgras{beretem}, \VUgras{sierżant sztabowy}
\textsuperscript{(27)} piechoty i \VUgras{sierżant artylerii} \textsuperscript{(28)} ze
złotym \VUgras{galonem} na rękawie, zatrzymała się blisko nas, żeby posłuchać wyjaśnień.

%%K t16-05-1 p
5. --- Mój dziadek, całkiem dumny z tego, że miał tak świetnych słuchaczy,
zaimprowizował kompletny plan bitwy.

%%K t16-05-2 p
Miasto ufortyfikowane, ze swoimi \VUgras{wałami} i \VUgras{fosami,} było oblegane przez
\VUgras{armię nieprzyjacielską,} która po długim \VUgras{bombardowaniu} była gotowa do
\VUgras{szturmu.} Ale \VUgras{oblężeni,} zmuszeni głodem, spróbowali \VUgras{wypadu} na
\VUgras{obóz} \VUgras{oblegających.} Odparłszy \VUgras{atak} zaciętą \VUgras{walką}
wokół \VUgras{obozowiska,} \VUgras{zwycięscy} oblegający puścili swoje
\VUgras{szwadrony,} żeby ścigać \VUgras{zwyciężonych} napastników. Ci \VUgras{cofnęli
się,} pokrywając \VUgras{pole bitwy} \VUgras{zabitymi} i \VUgras{rannymi.}
\VUgras{Noszowi} unieśli tych ostatnich do \VUgras{ambulansu,} żeby dać ich pod opiekę
\VUgras{chirurga.}

%%K t16-05-3 p
Mimo bohaterskiego oporu kilku \VUgras{batalionów,} które utworzyły \VUgras{czworobok,}
\VUgras{klęska} była zupełna, reszta armii \VUgras{uciekła,} zostawiając wielu
\VUgras{jeńców.} Nieprzyjaciel poszedł dopełnić swojego \VUgras{zwycięstwa,} zajmując
miasto. Może to będzie koniec kampanii i po \VUgras{rozejmie} będzie można podpisać
\VUgras{traktat pokojowy.}

%%K t16-tit-3 sub
\VUsaut{4.4mm}
\VUpk{10.2pt}{40}{Podarki na nowy rok.}
\VUsaut{2.8mm}

%%K t16-06-1 p
6. --- Młodzi żołnierze, którzy śmiali się, słuchając malowniczego opowiadania weterana,
poprosili go o kilka innych wyjaśnień. Co do mnie, obszedłem sklep, żeby popatrzeć na
piękną \VUgras{kuszę} \textsuperscript{(32)} i \VUgras{łuk} \textsuperscript{(33)} z jego
\VUgras{strzałą} \textsuperscript{(31)}. Tam znalazłem inny zbiór zwierząt: dwa
\VUgras{ptaki śpiewające} \textsuperscript{(34)}: automatycznego \VUgras{słowika} i
\VUgras{pokrzewkę} śpiewającą całym gardłem, i szereg dziwnych zwierząt:
\VUgras{nietoperza} \textsuperscript{(35)}, \VUgras{boa} \textsuperscript{(36)},
\VUgras{żółwia} \textsuperscript{(37)}, \VUgras{fokę} \textsuperscript{(38)},
\VUgras{wieloryba} \textsuperscript{(39)} i \VUgras{rekina} \textsuperscript{(40)} ze
złoconej skóry.

%%K t16-07-1 p
7. --- Nie zatrzymałem się długo, żeby je oglądać, bo dostrzegłem to, o czym marzyłem:
przepiękny \VUgras{teatr marionetek} \textsuperscript{(41)} ze wspaniałymi
\VUgras{dekoracjami} i zachwycającą czerwoną \VUgras{kurtyną.} Na \VUgras{scenie} dwaj
aktorzy zdawali się grać \VUgras{rolę} w bardzo ciekawej \VUgras{sztuce teatralnej,} bo
mali \VUgras{widzowie} z tektury umieszczeni w lożach albo siedzący na \VUgras{fotelach}
i na \VUgras{parterze} za \VUgras{dyrygentem} żywo bili brawo.

%%K t16-07-2 p
Niestety, ten teatr kosztował bardzo drogo.

%%K t16-08-1 p
8. --- Zresztą wybrać zabawki było trudno, bo w witrynach nagromadzono wszelkiego
rodzaju zabawki: \VUgras{Arlekina} \textsuperscript{(42)}, \VUgras{Pulcinellę}
\textsuperscript{(43)}, \VUgras{Pierrota} \textsuperscript{(44)}, \VUgras{puszczyki}
\textsuperscript{(45)} machające skrzydłami, gwiżdżące \VUgras{kosy}
\textsuperscript{(46)}, szczekające \VUgras{pudle}, \VUgras{fonograf}
\textsuperscript{(47)} i \VUgras{łamigłówki.} Jedna z tych zabawek bardzo mnie
rozbawiła: przedstawiała \VUgras{bitwę morską} \textsuperscript{(48)}, w której
\VUgras{statek} był atakowany przez \VUgras{piratów} blisko lądu obsadzonego
\VUgras{palmami kokosowymi.}

%%K t16-noto-1 noto
\VUnotes{143.2mm}{%
\textit{(*) Zobacz tablicę Nr 6.}%
}

%%K t16-09-1 p
9. --- Mogłem bardzo łatwo oglądać wszystkie te cuda, bo tylko mało osób było w
magazynie, ledwie kilku spacerowiczów, \VUgras{dragon} \textsuperscript{(49)} i
\VUgras{inkasent} \textsuperscript{(50)}, który przedstawiał \VUgras{tratę} w głównej
\VUgras{kasie} \textsuperscript{(51)}.

%%K t16-10-1 p
10. --- Zapytałem mojego dziadka, do czego służą książki położone na półkach za
\VUgras{kasjerką}; odpowiedział mi, że to \VUgras{księgi rachunkowe.}

%%K t16-tit-4 sub
\VUsaut{3.6mm}
\VUpk{10.2pt}{40}{Gapienie się po sklepie.}
\VUsaut{2.8mm}

%%K t16-11-1 p
11. --- Opuszczając dział zabawek, udaliśmy się do rymarni, żeby rzucić okiem na
\VUgras{siodła} \textsuperscript{(53)}, \VUgras{strzemiona} \textsuperscript{(54)},
\VUgras{uzdy} \textsuperscript{(55)}, \VUgras{wędzidła} \textsuperscript{(56)} i
\VUgras{szpicruty.} Podczas gdy \VUgras{kierownik działu} \textsuperscript{(57)} dawał
kilka informacji mojemu dziadkowi, poszedłem popatrzeć na wystawę \VUgras{porcelany i
kryształów} \textsuperscript{(58)}, gdzie są piękne \VUgras{salaterki} i delikatne
\VUgras{czajniki} \textsuperscript{(59)}.

%%K t16-12-1 p
12. --- Żeby pójść na pierwsze piętro, przeszliśmy za witryną \VUgras{gorsetów}
\textsuperscript{(60)}, obok działu pończoch, gdzie wystawiono bardzo piękne
\VUgras{pantalony} \textsuperscript{(63)}. Blisko \VUgras{sprzedawczyni}
\textsuperscript{(61)} kazała wybierać \VUgras{kupującej} \textsuperscript{(62)} piękne
jedwabne \VUgras{tkaniny} \textsuperscript{(64)}.

%%K t16-12-2 p
Wchodząc po schodach, zauważyłem, że magazyn jest ozdobiony japońskimi
\VUgras{lampionami} \textsuperscript{(65)}, \VUgras{parasolkami} \textsuperscript{(66)} i
ogromnymi \VUgras{palmami} \textsuperscript{(70)}.

%%K t16-13-1 p
13. --- Na pierwszym piętrze niewiele było do zobaczenia; tylko meble (*), \VUgras{kasy
pancerne} \textsuperscript{(67)}, \VUgras{komody} \textsuperscript{(68)},
\VUgras{wózki dziecięce} \textsuperscript{(69)}. Ale widok całego magazynu był bardzo
\VUgras{zabawny.}

%%K t16-14-1 p
14. --- U naszych stóp zobaczyłem instrumenty muzyczne: \VUgras{harfy}
\textsuperscript{(71)}, \VUgras{gitary} \textsuperscript{(72)}, \VUgras{mandoliny}
\textsuperscript{(73)}, \VUgras{kornety wentylowe} \textsuperscript{(74)},
\VUgras{klarnety} \textsuperscript{(75)}, skrzypce z ich \VUgras{smyczkiem}
\textsuperscript{(76)}, a nawet \VUgras{wielki bęben} \textsuperscript{(77)}. Nieco
dalej, przy dziale papierniczym, \VUgras{student} \textsuperscript{(78)} targował u
\VUgras{subiekta} \textsuperscript{(79)} teczkę na zeszyty. Blisko nich rozróżniłem
piękny \VUgras{elementarz} \textsuperscript{(80)}.

%%K t16-14-2 p
Pośrodku magazynu wzniesiono wysoką \VUgras{choinkę} \textsuperscript{(81)} całkiem
zaopatrzoną w świeczki, drobiazgi i wszelkiego rodzaju zabawki: \VUgras{drewniane konie}
\textsuperscript{(82)}, \VUgras{lalki} \textsuperscript{(83)}, itd..

%%K t16-14-3 p
\VUgras{Perfumeria} \textsuperscript{(84)} ze swoimi \VUgras{środkami do zębów} i
rozmaitymi \VUgras{perfumami} roztaczała bardzo dobrą woń, która dolatywała do nas.

%%K t16-tit-5 sub
\VUsaut{3.4mm}
\VUpk{10.2pt}{40}{Kupujemy coś.}
\VUsaut{2.8mm}

%%K t16-15-1 p
15. --- Ponieważ mój dziadek zaczął uznawać czas za zbyt długi, zeszliśmy wielkimi
schodami. Zatrzymał się nieco przed \VUgras{tablicą astronomiczną} \textsuperscript{(85)}
przedstawiającą, jak mi powiedział, \VUgras{komety} \textsuperscript{(a)},
\VUgras{zaćmienia księżyca i słońca} \textsuperscript{(b)}, różę wiatrów z
\VUgras{czterema stronami świata} \textsuperscript{(c)} \VUgras{(Północ, Południe,
Wschód i Zachód)}, mapę dwóch \VUgras{półkul} \textsuperscript{(d)},
\VUgras{planety} \textsuperscript{(e)} i \VUgras{układ słoneczny} \textsuperscript{(f)}.
Nie zrozumiałem nic z tego wszystkiego i o wiele bardziej zajęła mnie obszyta galonami
liberia \VUgras{posłańca} \textsuperscript{(86)}, który przechodził, i piękne
\VUgras{wachlarze} \textsuperscript{(87)}, które były obok mnie, blisko
\VUgras{portmonetek} \textsuperscript{(88)} i \VUgras{portfeli} \textsuperscript{(89)}.

%%K t16-16-1 p
16. --- Podziwiałem też na wystawie \VUgras{pióra} \textsuperscript{(90)} i
\VUgras{klamry do pasa} \textsuperscript{(91)}, piękne \VUgras{kwiaty sztuczne}
\textsuperscript{(94)} wdzięcznie ułożone w koszach. \VUgras{Fiołki},
\VUgras{lewkonie}, \VUgras{hiacynty} \textsuperscript{(95)}, \VUgras{pelargonie}
\textsuperscript{(96)}, \VUgras{lilie} \textsuperscript{(97)} i \VUgras{nasturcje}
\textsuperscript{(98)} były bardzo dobrze naśladowane.

%%K t16-tit-6 sub
\VUsaut{4.4mm}
\VUpk{10.2pt}{40}{Podarek mojego dziadka.}
\VUsaut{2.8mm}

%%K t16-17-1 p
17. --- Poprosiłem mojego dziadka, żeby kupił mi jedną ze \VUgras{szczoteczek do zębów}
\textsuperscript{(92)}, które były w pudełku obok \VUgras{szpilek do włosów}
\textsuperscript{(93)}. Zgodził się, a ja sam poszedłem zapłacić za swój zakup przy
kasie, blisko której przygotowywano ważną \VUgras{przesyłkę} \textsuperscript{(100)} dla
\VUgras{klienta} \textsuperscript{(99)}.

%%K t16-18-1 p
18. --- Nagle lekkie zamieszanie powstało wśród osób, które zatrzymały się blisko
artystycznych \VUgras{brązów} \textsuperscript{(101)}. \VUgras{Złodziej}
\textsuperscript{(102)} został właśnie schwytany na gorącym uczynku przez
\VUgras{inspektora} \textsuperscript{(103)}, gdy próbował kogoś okraść. Zadbano o to,
żeby sprowadzić policjanta, który miał go \VUgras{aresztować,} a właśnie kiedy
wychodziliśmy, sprawcę prowadzono do więzienia.
\VUsaut{6.0mm}
\VUfilet{22mm}
